\begin{proposition}[零温端的 Puiseux 精细展开：次增长分支的三项闭式]\label{prop:real-input-40-arity-charge-subbranch-puiseux}
\leanverified{paper\_real\_input\_40\_arity\_charge\_subbranch\_puiseux}
沿用定理 \ref{thm:real-input-40-arity-charge-det-closed} 的记号。除最大正实根分支 \(\lambda(q)=\rho(M(q))\) 外（其零温展开见命题 \ref{prop:real-input-40-arity-charge-zero-temp-expansion}），代数方程 \(P_7(\Lambda,q)=0\) 还存在一条正实增长分支 \(\lambda_{\mathrm{sub}}(q)\)。
当 \(q\to+\infty\) 时，该分支具有 Puiseux 渐近展开
\[
\boxed{\;
\lambda_{\mathrm{sub}}(q)=\sqrt{2q}+\frac{\sqrt2}{4}\,q^{-1/2}-\frac78\,q^{-1}+O(q^{-3/2})\; }.
\]
\end{proposition}
\begin{proof}
由定理 \ref{thm:real-input-40-arity-charge-det-closed}，\(\lambda_{\mathrm{sub}}(q)\) 为 \(P_7(\Lambda,q)=0\) 的某条正实代数分支。
置 \(q=t^{-2}\)（\(t\downarrow 0\)）并作 Puiseux 设形
\[
\lambda_{\mathrm{sub}}(q)=\frac{\sqrt2}{t}+b_1t+b_2t^2+O(t^3).
\]
将其代回 \(P_7(\lambda_{\mathrm{sub}}(t^{-2}),t^{-2})=0\)，并按 \(t\) 的幂次逐次比较系数，得
\(
b_1=\frac{\sqrt2}{4}
\)
与
\(
b_2=-\frac78
\)，
从而得到所示展开。证毕。
\end{proof}

\begin{proposition}[零荷端的解析起跳斜率：\texorpdfstring{$\lambda'(0)$}{lambda'(0)} 的闭式代数表达]\label{prop:real-input-40-arity-charge-derivative-q0}
\leanverified{paper\_real\_input\_40\_arity\_charge\_derivative\_q0}
沿用定理 \ref{thm:real-input-40-arity-charge-det-closed} 的记号。令 \(\kappa>1\) 为方程 \(\kappa^4-\kappa^3-1=0\) 的唯一实根（同命题 \ref{prop:real-input-40-arity-charge-zero-zeta}），则 \(\lambda(q)\) 在 \(q=0\) 的邻域解析，且满足
\[
\lambda(0)=\kappa,\qquad
\lambda(q)=\kappa+\lambda'(0)\,q+O(q^2),
\]
其中起跳斜率具有显式闭式
\[
\boxed{\;
\lambda'(0)=\frac{2\kappa^3+\kappa^2+8\kappa-2}{\kappa(4\kappa-3)}\; }.
\]
\end{proposition}
\begin{proof}
由定理 \ref{thm:real-input-40-arity-charge-det-closed} 中 \(P_7\) 的显式形式可直接分解
\[
P_7(\Lambda,0)=\Lambda^2(\Lambda-1)(\Lambda^4-\Lambda^3-1),
\]
从而 \(P_7(\kappa,0)=0\) 且 \(\kappa\) 为简单根。由隐函数定理，存在 \(q=0\) 邻域内的解析分支 \(\lambda(q)\) 满足 \(\lambda(0)=\kappa\) 与 \(P_7(\lambda(q),q)=0\)。
对 \(q\) 作隐式求导并在 \(q=0\) 处取值，得
\[
\lambda'(0)=-\frac{\partial_q P_7(\kappa,0)}{\partial_\Lambda P_7(\kappa,0)}.
\]
将 \(P_7\) 的显式系数代入并用 \(\kappa^4-\kappa^3-1=0\) 化简，即得所示有理闭式。证毕。
\end{proof}
